% ---------------------------------------------------------------
%  Preamble: default LaTeX fonts (Computer Modern), ISO 80000-2 style
% ---------------------------------------------------------------
\usepackage[b5paper,hmargin=19mm,top=20mm,bottom=22mm,headheight=16pt]{geometry}
\usepackage{emptypage}             % blank pages before odd-page chapter starts
\usepackage{amsmath,amssymb,amsthm,mathtools}
\usepackage{bm}
\usepackage[Euler]{upgreek}        % upright \uppi only (ISO: constants upright)
\usepackage{array,booktabs,multirow,tabularx}
\usepackage{enumitem}
\usepackage{xpatch}               % robust-command patches; loads etoolbox
\usepackage{xcolor}
\usepackage{tikz}
\usetikzlibrary{arrows.meta,calc,decorations.markings,patterns,positioning}
\usepackage{pgfplots}
\pgfplotsset{compat=1.18}
\usepackage[font=small,labelfont=md]{caption}
\usepackage{float}
\usepackage[most]{tcolorbox}
\usepackage{esvect}
\usepackage[colorlinks=true,linkcolor=blue!50!black,citecolor=blue!50!black,urlcolor=blue!60!black]{hyperref}
\usepackage[capitalise,nameinlink]{cleveref}

% Body text: 16 pt baseline spacing at the existing 11 pt class font size.
% Keep the separate small-text and heading sizes and spacing unchanged.
\xpatchcmd{\normalsize}{13.6}{16}{}%
  {\PackageError{lecture-notes}{Could not set 16 pt body line spacing}%
    {Check the 11 pt book class definition of \string\normalsize.}}
\normalsize

\setlist{itemsep=2pt,topsep=4pt}
\setlength{\parskip}{3pt}
\renewcommand{\arraystretch}{1.18}
\setlength{\emergencystretch}{2em}  % the B5 measure is narrow: allow looser lines, not overfull ones
\raggedbottom                       % no stretched vertical gaps to fill a page

% ---------------- page layout for B5 ----------------
% Section headings ragged right, so long titles break instead of overflowing.
\makeatletter
% Regular-weight text for tablet reading. Patch text styles individually so
% mathematical bold (\mat and \one) keeps its matrix/scalar distinction.
\newcommand{\regulartextstyle}[1]{%
  \patchcmd{#1}{\bfseries}{\mdseries}{}%
    {\PackageError{lecture-notes}{Could not set regular-weight text style}%
      {Check the book class definition of \string#1.}}}
\regulartextstyle{\@makechapterhead} % chapter number
\regulartextstyle{\@makechapterhead} % chapter title
\regulartextstyle{\@makeschapterhead}
\regulartextstyle{\l@chapter}
\regulartextstyle{\subsubsection}
\regulartextstyle{\paragraph}
\regulartextstyle{\subparagraph}
\regulartextstyle{\descriptionlabel}
\regulartextstyle{\labelitemii}
\thm@headfont{\mdseries}
\renewcommand\section{\@startsection{section}{1}{\z@}%
  {-3.5ex \@plus -1ex \@minus -.2ex}{2.3ex \@plus.2ex}%
  {\normalfont\Large\mdseries\raggedright}}
\renewcommand\subsection{\@startsection{subsection}{2}{\z@}%
  {-3.25ex\@plus -1ex \@minus -.2ex}{1.5ex \@plus .2ex}%
  {\normalfont\large\mdseries\raggedright}}
\makeatother
% Running heads in a smaller size, so long chapter and section titles fit the line.
\usepackage{fancyhdr}
\pagestyle{fancy}
\fancyhf{}
\fancyhead[LE,RO]{\thepage}
\fancyhead[RE]{\footnotesize\slshape\leftmark}
\fancyhead[LO]{\footnotesize\slshape\rightmark}
\renewcommand{\headrulewidth}{0pt}

% ---------------- ISO 80000-2 conventions ----------------
\newcommand{\dd}{\mathop{}\!\mathrm{d}}      % differential d (upright)
\newcommand{\ee}{\mathrm{e}}                  % Euler's number (upright)
\newcommand{\ii}{\mathrm{i}}                  % imaginary unit (upright)
\newcommand{\cpi}{\uppi}                      % the number pi (upright)
\DeclareRobustCommand{\vec}[1]{\vv{#1}}       % vectors: arrow over italic letter
\newcommand{\mat}[1]{\bm{#1}}                 % matrices: bold italic
\newcommand{\one}{\mat{1}}                    % unit matrix
\newcommand{\T}{^{\mathrm{T}}}                % transpose
\newcommand{\herm}{^{\dagger}}                % hermitian conjugate
\DeclareMathOperator{\tr}{tr}
\DeclareMathOperator{\diag}{diag}
\DeclareMathOperator{\sgn}{sgn}
\DeclareMathOperator{\Ind}{Ind}
\DeclareMathOperator{\Span}{span}
\DeclareMathOperator{\Ker}{Ker}
\DeclareMathOperator{\Imag}{Im}
\DeclareMathOperator{\Real}{Re}
\newcommand{\SO}[1]{\mathrm{SO}(#1)}
\newcommand{\SU}[1]{\mathrm{SU}(#1)}
\newcommand{\Og}[1]{\mathrm{O}(#1)}
\newcommand{\U}[1]{\mathrm{U}(#1)}
\newcommand{\GL}[1]{\mathrm{GL}(#1)}
% Schoenflies point-group names: letter italic, label subscript upright
\newcommand{\Cv}[1]{C_{#1\mathrm{v}}}
\newcommand{\Ch}[1]{C_{#1\mathrm{h}}}
\newcommand{\Dh}[1]{D_{#1\mathrm{h}}}
\newcommand{\Td}{T_{\mathrm{d}}}
\newcommand{\Oh}{O_{\mathrm{h}}}
\newcommand{\Ci}{C_{\mathrm{i}}}
\newcommand{\Cs}{C_{\mathrm{s}}}
% symmetry operations (group elements italic; labels upright)
\newcommand{\sv}{\sigma_{\mathrm{v}}}
\newcommand{\sd}{\sigma_{\mathrm{d}}}
\newcommand{\sh}{\sigma_{\mathrm{h}}}
\newcommand{\Ebar}{\bar{E}}
% irreducible representations (Mulliken labels): upright
\newcommand{\irr}[1]{\mathrm{#1}}
% representations with fixed meaning (upright labels)
\newcommand{\GV}{\Gamma_{\mathrm{V}}}
\newcommand{\GR}{\Gamma_{\mathrm{R}}}
\newcommand{\GM}{\Gamma_{\mathrm{M}}}
\newcommand{\Gvib}{\Gamma_{\mathrm{vib}}}
\newcommand{\chiV}{\chi_{\mathrm{V}}}
\newcommand{\chiR}{\chi_{\mathrm{R}}}
\newcommand{\chiM}{\chi_{\mathrm{M}}}
\newcommand{\Ninv}{N_{\mathrm{inv}}}
% misc
\DeclarePairedDelimiter{\abs}{\lvert}{\rvert}
\DeclarePairedDelimiter{\norm}{\lVert}{\rVert}
\newcommand{\ket}[1]{\lvert #1\rangle}
\newcommand{\bra}[1]{\langle #1\rvert}
\newcommand{\braket}[2]{\langle #1\vert #2\rangle}
\newcommand{\inner}[2]{\langle #1, #2\rangle}
\newcommand{\Tgroup}{\mathcal{T}}
\newcommand{\seitz}[2]{\{#1\,|\,#2\}}
\newcommand{\TR}{\Theta}

% ---------------- theorem-like environments ----------------
\theoremstyle{plain}
\newtheorem{theorem}{Theorem}[chapter]
\newtheorem{lemma}[theorem]{Lemma}
\newtheorem{proposition}[theorem]{Proposition}
\newtheorem{corollary}[theorem]{Corollary}
\theoremstyle{definition}
\newtheorem{definition}[theorem]{Definition}
\newtheorem{example}[theorem]{Example}
\theoremstyle{remark}
\newtheorem{remark}[theorem]{Remark}

% Box titles are white on a dark bar, so links inside a title must be white too.
\newcommand{\titlelinks}{\hypersetup{linkcolor=white,urlcolor=white}}
\tcbset{
  lecbox/.style={colback=blue!4,colframe=blue!45!black,boxrule=0.5pt,arc=1mm,
                 left=6pt,right=6pt,top=4pt,bottom=4pt,fonttitle=\mdseries\small\titlelinks},
  keybox/.style={colback=orange!6,colframe=orange!60!black,boxrule=0.5pt,arc=1mm,
                 left=6pt,right=6pt,top=4pt,bottom=4pt,fonttitle=\mdseries\small\titlelinks},
  gapbox/.style={colback=green!5,colframe=green!45!black,boxrule=0.5pt,arc=1mm,
                 left=6pt,right=6pt,top=4pt,bottom=4pt,fonttitle=\mdseries\small\titlelinks}
}
\newenvironment{watch}[1]{\begin{tcolorbox}[lecbox,title={Video: #1}]\small}{\end{tcolorbox}}
\newenvironment{keyresult}[1][Key result]{\begin{tcolorbox}[keybox,title={#1}]}{\end{tcolorbox}}
\newenvironment{beyond}[1][Beyond the lectures]{\begin{tcolorbox}[gapbox,title={#1}]\small}{\end{tcolorbox}}

% Sections and subsections that go beyond the lectures carry a green star.
\DeclareRobustCommand{\bstar}{\textcolor{green!45!black}{$\bigstar$}}
\newcommand{\beyondmark}{\texorpdfstring{\ \bstar}{ *}}

% ---------------- exercises and solutions ----------------
\newcommand{\exercisesection}{%
  \section*{Exercises and solutions}%
  \addcontentsline{toc}{section}{Exercises and solutions}%
  \markright{\MakeUppercase{Exercises and solutions}}}
\newcounter{exercise}[chapter]
\renewcommand{\theexercise}{\thechapter.\arabic{exercise}}
\newenvironment{exercise}[1][]{\refstepcounter{exercise}\par\medskip\noindent
  \textmd{Exercise~\theexercise}\if\relax\detokenize{#1}\relax\else\ \textit{(#1)}\fi.\ }{\par}
% Each solution follows its exercise, set off by a grey rule on the left.
\newtcolorbox{solution}{breakable,enhanced jigsaw,blanker,left=9pt,
  borderline west={1pt}{0pt}{black!35},before skip=4pt,after skip=10pt,
  before upper={\noindent\textit{Solution.}\ }}
\crefname{exercise}{Exercise}{Exercises}

% ---------------- links to the video lectures ----------------
\newcommand{\lecid}[1]{%
  \ifcase#1\or iHRC-f4SPxs\or 19uuuiJiY8Y\or BUZ-ZwHDl6M\or wTN-b-uhbAE\or
  XYyzpEt\string_Vgk\or vUt\string_IAdrrN8\or oZ5iq\string_\string_MQgk\or rXGrZq4mZAY\fi}
% \ts{lecture}{seconds}{mm:ss}: clickable timestamp (youtu.be short link)
\newcommand{\ts}[3]{\href{https://youtu.be/\lecid{#1}?t=#2}{L#1\,@\,#3}}
\newcommand{\lec}[1]{\href{https://youtu.be/\lecid{#1}}{Lecture~#1}}
